\section{Data and Compliance Posture}\label{data-and-compliance-posture}

\subsection{Source data}\label{source-data}

The classifier is fine-tuned on \textbf{MIMIC-IV-FHIR v2.1}
\autocite{Bennett2024,Bennett2023,Johnson2023} under PhysioNet
credentialed access. MIMIC-IV-FHIR is a FHIR R4 representation of
MIMIC-IV produced by the MIMIC team, distributed by PhysioNet under
their Credentialed Health Data License. We use three resource
types from this dataset:

\begin{itemize}
\tightlist
\item
  \texttt{MimicCondition.ndjson} --- diagnoses encoded as FHIR \texttt{Condition}
  resources with ICD-9-CM and ICD-10-CM codings under
  MIMIC-namespaced \texttt{CodeSystem} URIs.
\item
  \texttt{MimicPatient.ndjson} --- patient demographics for split-stratification
  metadata (not used as model input).
\item
  \texttt{MimicEncounter.ndjson} --- encounter context for joining conditions
  to encounters when needed by downstream code paths.
\end{itemize}

\subsection{Compliance posture}\label{compliance-posture}

PhysioNet's responsible-LLM-use policy \autocite{PhysioNet2025} explicitly
prohibits sharing credentialed data with third parties, including
``sending it through APIs.'' We adopt a \textbf{strictly local} posture:

\begin{itemize}
\tightlist
\item
  MIMIC content is downloaded directly from PhysioNet to the
  corresponding author's own laptop (Apple M1 MacBook Pro, 16 GB
  unified memory) via authenticated \texttt{wget}. Files live under a
  gitignored \texttt{data/mimic/raw/} directory.
\item
  The training pipeline reads parquet derivatives produced by a
  local preparation step. Both raw and derivative files remain on
  the laptop.
\item
  The training module sets defensive environment variables at
  import time to disable cloud telemetry libraries (wandb,
  mlflow, comet) so accidental imports cannot leak training metrics.
\item
  Trained checkpoints persist to a gitignored \texttt{models/checkpoints/}
  directory. Per PhysioNet policy, MIMIC-derivative weights are
  not redistributed via this repository, HuggingFace Hub, or any
  cloud bucket.
\item
  The evaluation runner that consumes the trained model reads
  predictions only from synthetic Synthea inputs (see
  \autocite{FungPhantomCodes2026}), so the trained classifier never produces
  predictions on MIMIC content that would then traverse a cloud
  endpoint.
\end{itemize}

The same code path supports CUDA and CPU backends; using either
would require an independent compliance review of the host
environment (e.g., Vertex AI Workbench, AWS EC2 with appropriate
BAAs). The local MPS path is the reference implementation because
it requires no such review.

\subsection{Cohort and label space}\label{cohort-and-label-space}

We restrict the cohort to ICD-10-CM codes in the CMS ACCESS Model
FHIR Implementation Guide v0.9.6 \autocite{CMS2026}, comprising two disease
groups:

\begin{itemize}
\tightlist
\item
  \textbf{CKM} (cardiometabolic): diabetes (E08/E09/E11/E13),
  atherosclerotic cardiovascular disease, chronic kidney disease
  stage 3.
\item
  \textbf{eCKM} (extended cardiometabolic): hypertension, dyslipidemia,
  prediabetes, obesity.
\end{itemize}

The MIMIC-IV-FHIR v2.1 source contains conditions coded under both
ICD-9-CM and ICD-10-CM, using MIMIC-namespaced FHIR \texttt{CodeSystem}
URIs (e.g.,
\texttt{http://mimic.mit.edu/fhir/mimic/CodeSystem/mimic-diagnosis-icd10})
rather than the canonical HL7 URIs, and using undotted code
formats (\texttt{E1110} rather than \texttt{E11.10}). We normalize both at the
extraction boundary: codes are converted to the canonical dotted
form (decimal point inserted after position 3 for codes longer
than three characters), and system URIs are mapped to the
canonical \texttt{http://hl7.org/fhir/sid/icd-10-cm} so that the rest of
the pipeline operates on a single canonical representation
regardless of source.

After ICD-10-CM filtering and ACCESS-scope filtering, the cohort
comprises \textbf{245,575 unique conditions} distributed 70/10/20
across train, validation, and test splits by stratified sampling
on \texttt{resource\_id} (ensuring all four degradation modes of a single
condition land in the same split). \textbf{178 unique ICD-10-CM codes}
appear in the train split. Each condition is materialized into
four rows --- one per degradation mode (D1\_full, D2\_no\_code,
D3\_text\_only, D4\_abbreviated) --- yielding \textbf{687,552 train rows,
98,272 validation rows, and 196,476 test rows}. The
classification head is sized for the \textbf{top-50} most frequent
codes, which together cover the head of the long-tail distribution
where supervised signal is densest.
